% Check criterion, deliverable, and the Decision Log question of Day 6.
% The criterion and the deliverable come from the Day 6 section of the
% syllabus and from Labs/day06/README.md.
\begin{frame}{Check criterion}
  \small
  The instructor checks this at the end of the lab:
  \begin{itemize}
    \item \cmark\ The notebook prints \code{complete} for tasks 1, 3, 4, 5,
          and 7.
    \item \cmark\ The report has the two predictions: task 2 and task 6.
          You wrote them before the experiments.
    \item \cmark\ The plot \code{tradeoff.png} has at least six models.
    \item \cmark\ The report names one model for the XIAO ESP32S3 and one
          model for the Raspberry Pi 5, with the numbers of each model.
    \item \cmark\ The selection states the constraint that decides for
          each product.
    \item \cmark\ The Decision Log gives numbers and names one trade-off.
  \end{itemize}
  \vskip 0.4em
  \textbf{Hand in:} the file \code{report.md} and the plot
  \code{tradeoff.png}.
\end{frame}

\begin{frame}{Decision Log}
  Write about 100 words. State one design decision, give your measured
  numbers, and name the trade-off.
  \vskip 0.4em
  \begin{exerciseblock}
    Which model do you select for the XIAO ESP32S3, and which model for the
    Raspberry Pi 5? Use the budgets of section 10 of the notebook.
  \end{exerciseblock}
  \vskip 0.2em
  A good Decision Log:
  \begin{itemize}
    \item gives the file size, the peak, the accuracy, and the latency of
          each selected model,
    \item names the constraint that decides for each product, and the
          model that this constraint removes,
    \item says what the product loses with the selection. Example: some
          points of accuracy against half of the RAM.
  \end{itemize}
\end{frame}
